\newpage
\section{软件编程}

\subsection{编程语言}

\subsubsection{Python}
Python是一种高级编程语言，以其简单易读的语法和广泛的库支持而闻名。
适用于数据分析、机器学习、Web开发等领域。

\subsubsection{Ruby}
Ruby是一种动态、面向对象的编程语言，以简洁和优雅的代码风格著称。
它常用于Web开发，特别是Ruby on Rails框架。

\subsubsection{Bash\cite{gnuBashManual}}
Bash是一种Unix shell和命令语言，用于编写脚本以自动化系统任务。
它是Linux和macOS中的默认shell。

\subsubsection{Lua}
Lua是一种轻量级的脚本语言，设计用于嵌入应用程序。
以其简单和高效著称，常用于游戏开发。

\subsubsection{C}
C是一种通用的编程语言，以其高效和灵活性广泛应用于系统编程、嵌入式系统和操作系统开发。

\subsubsection{C++}
C++是C语言的扩展，支持面向对象编程。
它用于开发高性能应用程序，如游戏、图形处理和大型软件系统。

\subsubsection{Java}
Java是一种广泛使用的编程语言，具有平台独立性和丰富的库支持。
常用于企业级应用、Android开发和Web服务。

\subsubsection{Kotlin}
Kotlin是一种现代编程语言，与Java高度兼容。
它简化了Android开发，并支持多平台应用程序开发。

\subsubsection{Swift}
Swift是苹果公司开发的编程语言，专为iOS和macOS应用程序设计。
以其安全性和高效性能而受到开发者欢迎。

\subsection{文档语言}

\subsubsection{wikitex}
MediaWiki使用的标记语言，适用于创建和编辑维基百科等Wiki网站。
它支持链接、表格、列表、模板等多种格式，便于协作编辑和内容管理。

\subsubsection{markdown}
Markdown是一种轻量化的文字排版标记语言，旨在以简单易读的纯文本格式编写文档。
它支持基本的文本格式，如标题、列表、链接、图片等，并可通过扩展支持复杂功能。
Markdown广泛用于博客、论坛和文档编写。

\subsubsection{latex}
LaTeX是一种用于复杂排版和数学公式的标记语言，特别适合学术论文、技术文档和书籍的编写。
它提供了强大的排版功能和宏包支持，使用户能够轻松处理跨页布局、索引、参考文献等。

\subsubsection{bibtex}
BibTeX是一种用于管理和格式化参考文献的标记语言，通常与LaTeX结合使用。
用户通过.bib文件管理文献条目，并在LaTeX文档中引用这些条目。
BibTeX支持多种文献类型，如书籍、文章、会议论文等，并提供多种引用样式，如APA、IEEE等。